% =====================================================================
% Project / report template for NUS Course Work
% Copy this file into a course's projects/ folder and rename it,
% e.g. CS2040/projects/assignment-1.tex
% Compile with: pdflatex assignment-1.tex
% (If you use the minted package for code, compile with:
%  pdflatex -shell-escape assignment-1.tex)
% =====================================================================
\documentclass[11pt,a4paper]{article}

\usepackage[utf8]{inputenc}
\usepackage[T1]{fontenc}
\usepackage[margin=1in]{geometry}
\usepackage{amsmath, amssymb}
\usepackage{graphicx}
\usepackage{enumitem}
\usepackage{listings}      % lightweight code listings (no shell-escape needed)
\usepackage{xcolor}
\usepackage{hyperref}
\hypersetup{colorlinks=true, linkcolor=blue, urlcolor=blue}

% --- Code listing style ---
\lstset{
    basicstyle=\ttfamily\small,
    keywordstyle=\color{blue},
    commentstyle=\color{gray},
    stringstyle=\color{red!60!black},
    numbers=left,
    numberstyle=\tiny\color{gray},
    breaklines=true,
    frame=single,
    showstringspaces=false
}

% --- Document metadata ---
\title{Course Code --- Project / Assignment Title}
\author{Your Name \\ Student No. AXXXXXXXX}
\date{\today}

\begin{document}
\maketitle

\begin{abstract}
A short summary of the project: the problem, your approach, and the result.
\end{abstract}

\tableofcontents
\bigskip

\section{Introduction}
Describe the problem and objectives.

\section{Approach}
Explain your method or design.

\section{Implementation}
Discuss the implementation. Example code listing:
\begin{lstlisting}[language=Python]
def hello():
    print("Hello, NUS!")
\end{lstlisting}

\section{Results}
Present results, figures, and analysis.

\section{Conclusion}
Summarize findings and possible improvements.

\end{document}
